\ssr{РЕФЕРАТ}

Расчетно-пояснительная записка \pageref{LastPage} страниц, \totalfigures{} рисунков, \totaltables{} таблиц, 45 источников, 3 приложения.

Ключевые слова: КОМПЬЮТЕРНАЯ ЛИНГВИСТИКА, СИНТАКСИЧЕСКИЙ АНАЛИЗ, ОГРАНИЧЕННЫЙ ЕСТЕСТВЕННЫЙ ЯЗЫК, ДЕРЕВО СОСТАВЛЯЮЩИХ, РУССКИЙ ЯЗЫК, АЛГОРИТМ КОКА–ЯНГЕРА–КАСАМИ.

Цель работы~---~разработка метода построения синтаксических деревьев для ограниченного естественного русского языка на основе наиболее употребимых грамматик.


В данной работе выполнен анализ предметной области автоматизированного построения синтаксических деревьев ограниченного естественного русского языка, проведен сравнительный анализ существующих методов построения синтаксических деревьев. Разработаны метод построения синтаксических деревьев и программное обеспечение реализующее описанный метод. Проведено исследование зависимости времени построения синтаксических деревьев от количества словоформ в предложении, приведена оценка точности и полноты разработанного метода.